% =================================================
% TikZ-рисунки документа
%
% Геометрия каждого рисунка хранится здесь один раз.
% Цвета и линии задаются в style-color.sty или style-bw.sty.
% =================================================

% Библиотеки для новых рисунков (серединный перпендикуляр, биссектриса):
% угловые дуги \pic{angle=...} и \pic{right angle=...}.
% calc уже подключён в style-common.sty.
\usetikzlibrary{angles,quotes}


% =================================================
% Общие стили для рисунков про серединный перпендикуляр
% и биссектрису угла. Префикс gmtpb отделяет их от стилей
% gmt * в style-color.sty / style-bw.sty.
% =================================================
\tikzset{
	gmtpb point/.style={
		circle,
		fill=black,
		inner sep=1.5pt
	},
	gmtpb main/.style={
		line width=0.8pt
	},
	gmtpb auxiliary/.style={
		line width=0.65pt,
		dashed,
		gray
	},
	gmtpb construction/.style={
		line width=0.65pt
	},
	gmtpb label/.style={
		font=\small
	},
	gmtpb caption/.style={
		font=\small,
		align=center,
		text width=5cm
	}
}

% =================================================
% Общие стили для рисунков «расстояние от точки до прямой»
% и «пересечение двух окружностей». Префикс gmtd отделяет
% их от стилей gmt * и gmtpb *.
% =================================================
\tikzset{
	gmtd point/.style={
		circle,
		fill=black,
		inner sep=1.3pt
	},
	gmtd main line/.style={
		line width=0.8pt
	},
	gmtd auxiliary/.style={
		draw=black!45,
		line width=0.65pt,
		dashed
	},
	gmtd diagram title/.style={
		font=\small\bfseries,
		align=center
	},
	gmtd diagram caption/.style={
		font=\small,
		align=center,
		text width=5.4cm
	}
}


% =================================================
% Рисунок 1.
% От нескольких подходящих точек к окружности
% =================================================

\newcommand{\gmtcirclefigure}{%
	\begin{center}
		\begin{tikzpicture}[
			scale=0.95,
			transform shape,
			every node/.style={
				font=\footnotesize
			}
			]
			
			\def\radius{2.2}
			
			% -----------------------------------------
			% Левая часть
			% -----------------------------------------
			\begin{scope}[shift={(0,0)}]
				
				\node[
				gmt heading,
				font=\small\bfseries\sffamily,
				text width=5.2cm,
				align=center
				] at (0,3.2)
				{Отмечены несколько подходящих точек};
				
				\draw[gmt auxiliary]
				(0,0) circle (\radius);
				
				\node[
				gmt center,
				label={
					[gmt center label]
					below left:$O$
				}
				] (O1) at (0,0) {};
				
				\coordinate (M1) at (25:\radius);
				\coordinate (M2) at (105:\radius);
				\coordinate (M3) at (195:\radius);
				\coordinate (M4) at (295:\radius);
				
				\draw[gmt radius] (O1) -- (M1);
				\draw[gmt radius] (O1) -- (M2);
				\draw[gmt radius] (O1) -- (M3);
				\draw[gmt radius] (O1) -- (M4);
				
				\node[
				gmt point,
				label={
					[gmt point label]
					above right:$M_1$
				}
				] at (M1) {};
				
				\node[
				gmt point,
				label={
					[gmt point label]
					above left:$M_2$
				}
				] at (M2) {};
				
				\node[
				gmt point,
				label={
					[gmt point label]
					below left:$M_3$
				}
				] at (M3) {};
				
				\node[
				gmt point,
				label={
					[gmt point label]
					below right:$M_4$
				}
				] at (M4) {};
				
				\node[
				above,
				gmt radius label
				] at ($(O1)!0.5!(M1)$)
				{$R$};
				
				\node[
				gmt caption,
				text width=6cm,
				align=center
				] at (0,-3.25)
				{Все отмеченные точки удовлетворяют\\[2pt]
					условию \(OM=R\)};
				
			\end{scope}
			
			% -----------------------------------------
			% Переход
			% -----------------------------------------
			\draw[gmt arrow]
			(3.15,-1) -- (4.15,-1);
			
			\node[
			gmt transition text,
			text width=1.8cm,
			align=center
			] at (3.65,0.65)
			{отметим\\все точки};
			
			% -----------------------------------------
			% Правая часть
			% -----------------------------------------
			\begin{scope}[shift={(7.3,0)}]
				
				\node[
				gmt heading,
				font=\small\bfseries\sffamily,
				text width=4.8cm,
				align=center
				] at (0,3.2)
				{Получается окружность};
				
				\draw[gmt main circle]
				(0,0) circle (\radius);
				
				\node[
				gmt center,
				label={
					[gmt center label]
					below left:$O$
				}
				] (O2) at (0,0) {};
				
				\coordinate (M) at (35:\radius);
				
				\draw[gmt radius]
				(O2) -- (M);
				
				\node[
				gmt point,
				label={
					[gmt point label]
					above right:$M$
				}
				] at (M) {};
				
				\node[
				above,
				gmt radius label
				] at ($(O2)!0.5!(M)$)
				{$R$};
				
				\node[
				gmt caption,
				text width=6cm,
				align=center
				] at (0,-3.25)
				{Все точки \(M\), для которых \(OM=R\),\\[2pt]
					образуют окружность};
				
			\end{scope}
			
		\end{tikzpicture}
	\end{center}
}


% =================================================
% Рисунок 2.
% Сравнение условий OM<R, OM=R, OM<=R, OM>R
% =================================================

\newcommand{\gmtdistancecomparisonfigure}{%
	\begin{center}
		\begin{tikzpicture}[
			scale=0.95,
			transform shape,
			every node/.style={
				font=\footnotesize
			}
			]
			
			% -----------------------------------------
			% Основные параметры
			% -----------------------------------------
			\def\radius{1.25}
			\def\panelshift{3.7}
			
			% -----------------------------------------
			% 1. OM < R
			% -----------------------------------------
			\begin{scope}[shift={(0,0)}]
				
				\node[
				gmt distance heading
				] at (0,2.05)
				{\(OM<R\)};
				
				% Внутренняя область входит,
				% окружность не входит.
				\fill[
				gmt distance inside fill
				]
				(0,0) circle (\radius);
				
				\draw[
				gmt distance excluded boundary
				]
				(0,0) circle (\radius);
				
				\node[
				gmt distance center,
				label={
					[gmt distance center label]
					below left:$O$
				}
				] (O1) at (0,0) {};
				
			\coordinate (M1) at (35:0.55);

			\draw[
			gmt distance radius
			]
			(O1) -- (M1);

			\node[
			gmt distance point,
			label={
				[gmt distance point label]
				right:$M$
			}
			] at (M1) {};
				
				\node[
				gmt distance caption,
				text width=2.8cm,
				align=center
				] at (0,-1.75)
				{внутренняя область\\без границы};
				
			\end{scope}
			
			% -----------------------------------------
			% 2. OM = R
			% -----------------------------------------
			\begin{scope}[shift={(\panelshift,0)}]
				
				\node[
				gmt distance heading
				] at (0,2.05)
				{\(OM=R\)};
				
				% Входит только окружность.
				\draw[
				gmt distance boundary,
				line width=1.1pt
				]
				(0,0) circle (\radius);
				
				\node[
				gmt distance center,
				label={
					[gmt distance center label]
					below left:$O$
				}
				] (O2) at (0,0) {};
				
			\coordinate (M2) at (35:\radius);

			\draw[
			gmt distance radius
			]
			(O2) -- (M2);

			\node[
			gmt distance point,
			label={
				[gmt distance point label]
				right:$M$
			}
			] at (M2) {};
				
				\node[
				gmt distance caption,
				text width=2.8cm,
				align=center
				] at (0,-1.75)
				{только\\окружность};
				
			\end{scope}
			
			% -----------------------------------------
			% 3. OM <= R
			% -----------------------------------------
			\begin{scope}[shift={(2*\panelshift,0)}]
				
				\node[
				gmt distance heading
				] at (0,2.05)
				{\(OM\leq R\)};
				
				% Входят внутренняя область
				% и окружность.
				\fill[
				gmt distance inside fill
				]
				(0,0) circle (\radius);
				
				\draw[
				gmt distance boundary,
				line width=1.1pt
				]
				(0,0) circle (\radius);
				
				\node[
				gmt distance center,
				label={
					[gmt distance center label]
					below left:$O$
				}
				] (O3) at (0,0) {};
				
			\coordinate (M3) at (35:0.72);

			\draw[
			gmt distance radius
			]
			(O3) -- (M3);

			\node[
			gmt distance point,
			label={
				[gmt distance point label]
				right:$M$
			}
			] at (M3) {};
				
				\node[
				gmt distance caption,
				text width=2.8cm,
				align=center
				] at (0,-1.75)
				{круг вместе\\с границей};
				
			\end{scope}
			
			% -----------------------------------------
			% 4. OM > R
			% -----------------------------------------
			\begin{scope}[shift={(3*\panelshift,0)}]
				
				\node[
				gmt distance heading
				] at (0,2.05)
				{\(OM>R\)};
				
				% Показывается ограниченная часть
				% внешней области.
				\fill[
				gmt distance outside fill,
				even odd rule
				]
				(-1.65,-1.55)
				rectangle
				(1.65,1.55)
				
				(0,0)
				circle
				(\radius);
				
				% Окружность не входит.
				\draw[
				gmt distance excluded boundary
				]
				(0,0) circle (\radius);
				
				\node[
				gmt distance center,
				label={
					[gmt distance center label]
					below left:$O$
				}
				] (O4) at (0,0) {};
				
			\coordinate (M4) at (25:1.48);

			\draw[
			gmt distance radius
			]
			(O4) -- (M4);

			\node[
			gmt distance point,
			label={
				[gmt distance point label]
				right:$M$
			}
			] at (M4) {};
				
				\node[
				gmt distance caption,
				text width=2.8cm,
				align=center
				] at (0,-1.75)
				{внешняя область\\без границы};
				
			\end{scope}
			
		\end{tikzpicture}
	\end{center}
}

% =================================================
% Рисунок 3.
% Полуокружность, круг и окружность:
% пропущенные и лишние точки
% =================================================

\newcommand{\gmtextraandmissingpointsfigure}{%
	\begin{center}
		\begin{tikzpicture}

			\def\radius{1.15}

			% =================================================
			% 1. Полуокружность
			% =================================================

			\begin{scope}[shift={(-3.1,0)}]

				\node[gmt error heading]
					at (0,2.15)
					{Полуокружность};

				\draw[gmt helper boundary]
					(0,0) circle (\radius);

				\draw[
					gmt correct boundary,
					line width=1.1pt
				]
					(-\radius,0)
					arc[
						start angle=180,
						end angle=0,
						radius=\radius
					];

				\node[
					gmt error center,
					label={
						[gmt error center label]
						below left:$O$
					}
				] at (0,0) {};

				\node[gmt correct point]
					at (140:\radius) {};

				\node[gmt correct point]
					at (90:\radius) {};

				\node[gmt correct point]
					at (40:\radius) {};

				\node[gmt highlighted error point]
					at (220:\radius) {};

				\node[gmt highlighted error point]
					at (270:\radius) {};

				\node[gmt highlighted error point]
					at (320:\radius) {};

				\node[
					gmt error callout text,
					align=center
				] at (0,-2.15)
					{пропущенные точки};

				\node[
					gmt error caption,
					align=center,
					text width=4.2cm
				] at (0,-3.25)
					{Пропущены точки, которые тоже
					удовлетворяют условию \(OM=R\).};

			\end{scope}

			% =================================================
			% 2. Круг
			% =================================================

			\begin{scope}[shift={(3.1,0)}]

				\node[gmt error heading]
					at (0,2.15)
					{Круг};

				\fill[gmt extra region]
					(0,0) circle (\radius);

				\draw[
					gmt correct boundary,
					line width=1.1pt
				]
					(0,0) circle (\radius);

				\node[
					gmt error center,
					label={
						[gmt error center label]
						below left:$O$
					}
				] at (0,0) {};

				\node[gmt correct point]
					at (135:\radius) {};

				\node[gmt correct point]
					at (40:\radius) {};

				\node[gmt highlighted error point]
					at (25:0.58) {};

				\node[gmt highlighted error point]
					at (155:0.68) {};

				\node[gmt highlighted error point]
					at (260:0.52) {};

				\node[
					gmt error callout text,
					align=center
				] at (0,-2.15)
					{лишние точки};

				\node[
					gmt error caption,
					align=center,
					text width=4.2cm
				] at (0,-3.25)
					{Добавлены лишние точки,
					для которых \(OM<R\).};

			\end{scope}

			% =================================================
			% 3. Окружность
			% =================================================

			\begin{scope}[shift={(0,-7.7)}]

				\node[gmt success heading]
					at (0,2.15)
					{Окружность};

				\draw[
					gmt correct boundary,
					line width=1.1pt
				]
					(0,0) circle (\radius);

				\node[
					gmt error center,
					label={
						[gmt error center label]
						below left:$O$
					}
				] (O) at (0,0) {};

				\coordinate (M)
					at (35:\radius);

				\draw[gmt error radius]
					(O) -- (M);

				\node[
					gmt correct point,
					label={
						[gmt error point label]
						above right:$M$
					}
				] at (M) {};

				\node[gmt correct point]
					at (140:\radius) {};

				\node[gmt correct point]
					at (250:\radius) {};

				\node[
					gmt success caption,
					align=center,
					text width=4.4cm
				] at (0,-2.25)
					{Нет лишних и пропущенных точек.\\
					\textbf{Это искомое ГМТ.}};

			\end{scope}

		\end{tikzpicture}
	\end{center}
}


% =================================================
% Рисунок 4.
% Серединный перпендикуляр к отрезку
% =================================================

\newcommand{\gmtperpendicularbisectorfigure}{%
	\begin{center}
		\begin{tikzpicture}[
			line cap=round,
			line join=round
		]

			% -----------------------------------------
			% Отрезок AB и серединный перпендикуляр
			% -----------------------------------------
			\begin{scope}[shift={(0,0)}]

				\coordinate (A) at (-2.4,0);
				\coordinate (B) at ( 2.4,0);
				\coordinate (K) at ( 0,0);
				\coordinate (M) at ( 0,3.1);

				% Отрезок AB
				\draw[gmtpb main] (A)--(B);

				% Серединный перпендикуляр
				\draw[gmtpb main] (0,-0.75)--(0,3.45);

				% Равные отрезки MA и MB
				\draw[gmtpb main] (M)--(A);
				\draw[gmtpb main] (M)--(B);

				% Засечки на AK и KB
				\draw[gmtpb construction]
					($(A)!0.5!(K)+(0,0.12)$)
					--
					($(A)!0.5!(K)+(0,-0.12)$);
				\draw[gmtpb construction]
					($(K)!0.5!(B)+(0,0.12)$)
					--
					($(K)!0.5!(B)+(0,-0.12)$);

				% Прямой угол
				\draw[gmtpb construction]
					($(K)+(0,0.24)$)
					--
					($(K)+(0.24,0.24)$)
					--
					($(K)+(0.24,0)$);

				% Точки и подписи
				\node[gmtpb point] at (A) {};
				\node[gmtpb point] at (B) {};
				\node[gmtpb point] at (K) {};
				\node[gmtpb point] at (M) {};
				\node[gmtpb label,below left=2pt]  at (A) {$A$};
				\node[gmtpb label,below right=2pt] at (B) {$B$};
				\node[gmtpb label,below right=3pt] at (K) {$K$};
				\node[gmtpb label,above left=2pt]  at (M) {$M$};

				% Подпись
				\node[gmtpb caption] at (0,-1.25)
					{Серединный перпендикуляр к отрезку $AB$};

			\end{scope}

		\end{tikzpicture}
	\end{center}
}


% =================================================
% Рисунок 5.
% Биссектриса угла
% =================================================

\newcommand{\gmtbisectorfigure}{%
	\begin{center}
		\begin{tikzpicture}[
			line cap=round,
			line join=round
		]

			% -----------------------------------------
			% Биссектриса угла и равные расстояния
			% -----------------------------------------
			\begin{scope}[shift={(0,0)}]

				\coordinate (O) at (0,0);
				\coordinate (A) at (4.4,0);
				\coordinate (B) at (65:4.15);

				% Биссектриса угла
				\coordinate (M) at (32.5:2.85);

				% Основания перпендикуляров
				\coordinate (P) at ($(O)!(M)!(A)$);
				\coordinate (Q) at ($(O)!(M)!(B)$);

				% Стороны угла
				\draw[gmtpb main] (O)--(A);
				\draw[gmtpb main] (O)--(B);

				% Биссектриса
				\draw[gmtpb main] (O)--(M);

				% Перпендикуляры к сторонам
				\draw[gmtpb auxiliary] (M)--(P);
				\draw[gmtpb auxiliary] (M)--(Q);

				% Дуги равных углов
				\pic[
					draw,
					angle radius=8mm,
					angle eccentricity=1.25
				] {angle=A--O--M};

				\pic[
					draw,
					angle radius=8mm,
					angle eccentricity=1.25
				] {angle=M--O--B};

				% Прямые углы
				\pic[
					draw,
					angle radius=3.5mm
				] {right angle=M--P--O};

				\pic[
					draw,
					angle radius=3.5mm
				] {right angle=O--Q--M};

				% Засечки на MP и MQ
				\coordinate (T_P) at ($(M)!0.5!(P)$);
				\coordinate (T_Q) at ($(M)!0.5!(Q)$);
				\draw[gmtpb construction]
					($(T_P)+(-0.10,0)$)
					--
					($(T_P)+(0.10,0)$);
				\draw[gmtpb construction,rotate around={65:(T_Q)}]
					($(T_Q)+(-0.10,0)$)
					--
					($(T_Q)+(0.10,0)$);

				% Точки
				\node[gmtpb point] at (O) {};
				\node[gmtpb point] at (A) {};
				\node[gmtpb point] at (B) {};
				\node[gmtpb point] at (M) {};
				\node[gmtpb point] at (P) {};
				\node[gmtpb point] at (Q) {};

				% Подписи
				\node[gmtpb label,below left=2pt]  at (O) {$O$};
				\node[gmtpb label,below=3pt]       at (A) {$A$};
				\node[gmtpb label,above=3pt]       at (B) {$B$};
				\node[gmtpb label,above right=2pt] at (M) {$M$};
				\node[gmtpb label,below=3pt]       at (P) {$P$};
				\node[gmtpb label,above right=1pt] at (Q) {$Q$};

				% Подпись
				\node[gmtpb caption] at (2.15,-1.25)
					{Биссектриса угла и расстояния от точки до его сторон};

			\end{scope}

		\end{tikzpicture}
	\end{center}
}


% =================================================
% Рисунок 6.
% Расстояние от точки до прямой
% =================================================

\newcommand{\gmtdistancetolinefigure}{%
	\begin{center}
		\begin{tikzpicture}[
			font=\small
		]

			% -----------------------------------------
			% Точка, прямая, перпендикуляр MH
			% -----------------------------------------
			\begin{scope}[shift={(0,0)}]

				\node[gmtd diagram title] at (0,3.05)
					{Расстояние от точки до прямой};

				\coordinate (M) at (0,2.05);
				\coordinate (H) at (0,0);

				% Прямая a
				\draw[gmtd main line]
					(-2.6,0) -- (2.6,0)
					node[right] {$a$};

				% Перпендикуляр MH
				\draw[gmtd main line]
					(M) -- (H);

				% Знак прямого угла
				\draw[line width=0.6pt]
					($(H)+(-0.18,0)$)
					-- ($(H)+(-0.18,0.18)$)
					-- ($(H)+(0,0.18)$);

				% Точки и подписи
				\node[gmtd point] at (M) {};
				\node[above right=1pt] at (M) {$M$};
				\node[gmtd point] at (H) {};
				\node[below=3pt] at (H) {$H$};

				% Подпись перпендикулярности
				\node[
					anchor=west,
					fill=white,
					inner sep=1.5pt
				] at ($(M)!0.52!(H)+(0.08,0)$)
					{$MH\perp a$};

				% Пояснение
				\node[gmtd diagram caption] at (0,-1.45)
					{Расстояние от точки \(M\) до прямой \(a\)
					 равно длине отрезка перпендикуляра \(MH\).};

			\end{scope}

		\end{tikzpicture}
	\end{center}
}


% =================================================
% Рисунок 7.
% Пересечение двух окружностей
% =================================================

\newcommand{\gmtcirclesintersectionfigure}{%
	\begin{center}
		\begin{tikzpicture}[
			font=\small
		]

			% -----------------------------------------
			% Две окружности и их общие точки
			% -----------------------------------------
			\begin{scope}[shift={(0,0)}]

				\node[gmtd diagram title] at (0,3.05)
					{Пересечение двух окружностей};

				% Центры
				\coordinate (A) at (-1.45,0);
				\coordinate (B) at ( 1.45,0);
				\def\radiusA{2.05}
				\def\radiusB{1.55}
				\def\centerdistance{2.9}

				% Точки пересечения вычисляются из радиусов окружностей,
				% поэтому маркеры всегда точно лежат на обеих окружностях.
				\pgfmathsetmacro{\intersectionoffset}{
					(\radiusA*\radiusA-\radiusB*\radiusB
					 +\centerdistance*\centerdistance)/(2*\centerdistance)
				}
				\pgfmathsetmacro{\intersectionheight}{
					sqrt(\radiusA*\radiusA
					 -\intersectionoffset*\intersectionoffset)
				}
				\coordinate (M1) at
					($(A)+(\intersectionoffset,\intersectionheight)$);
				\coordinate (M2) at
					($(A)+(\intersectionoffset,-\intersectionheight)$);

				% Окружности
				\draw[gmtd main line]
					(A) circle (\radiusA);
				\draw[gmtd main line]
					(B) circle (\radiusB);

				% Линия центров
				\draw[gmtd auxiliary]
					(A) -- (B);

				% Радиусы к одной точке пересечения
				\draw[gmtd main line]
					(A) -- (M1);
				\draw[gmtd main line]
					(B) -- (M1);

				% Центры
				\node[gmtd point] at (A) {};
				\node[below left=2pt] at (A) {$A$};
				\node[gmtd point] at (B) {};
				\node[below right=2pt] at (B) {$B$};

				% Точки пересечения
				\node[gmtd point] at (M1) {};
				\node[above=3pt] at (M1) {$M_1$};
				\node[gmtd point] at (M2) {};
				\node[below=3pt] at (M2) {$M_2$};

				% Пояснение
				\node[gmtd diagram caption] at (0,-2.55)
					{Точки \(M_1\) и \(M_2\) принадлежат обеим окружностям.};

			\end{scope}

		\end{tikzpicture}
	\end{center}
}
